\chapter{欧氏空间中的微分形式}
\label{chap:euclidean-differential-forms-solutions}

\begin{solution}{alternating-form-basis}
交错 $k$-线性形式由其在标准基向量组上的值决定；若指标重复则值为零，若
指标互异则只差一个置换符号。因此
$dx^{i_1}\wedge\cdots\wedge dx^{i_k}$（$i_1<\cdots<i_k$）张成且线性无关，
共有 $\binom nk$ 个。$n=4,k=2$ 时标准基为
\[
 dx^1\wedge dx^2,dx^1\wedge dx^3,dx^1\wedge dx^4,
 dx^2\wedge dx^3,dx^2\wedge dx^4,dx^3\wedge dx^4.
\]
\end{solution}

\begin{solution}{wedge-sign-computation}
例如
$(dx+dy)\wedge(dx-dy)=-2dx\wedge dy$，因为
$dx\wedge dx=dy\wedge dy=0$ 且 $dy\wedge dx=-dx\wedge dy$。
一般地，一项若含两个相同的一次形式，交换这两个因子一方面不改变该项，
另一方面产生负号，故该项等于自身的负数，只能为零。
\end{solution}

\begin{solution}{top-form-determinant-pullback}
设线性映射矩阵为 $A=(a_{ij})$。有
$A^*dx^i=\sum_j a_{ij}dx^j$。展开
\[
 A^*(dx^1\wedge\cdots\wedge dx^n)
 =\bigwedge_{i=1}^n\sum_ja_{ij}dx^j.
\]
含重复指标的项全为零，剩余项按置换求和，其系数正是 $\det A$。因此
\[
 A^*(dx^1\wedge\cdots\wedge dx^n)
 =(\det A)\,dx^1\wedge\cdots\wedge dx^n.
\]
\end{solution}

\begin{solution}{exterior-derivative-leibniz}
对齐次形式 $\alpha=\sum_Ia_I dx^I$、$\beta=\sum_Jb_Jdx^J$，用普通乘积法则
$d(a_Ib_J)=b_Jda_I+a_Idb_J$。把 $db_J$ 从 $dx^I$ 左侧移到右侧要交换
$\deg\alpha$ 次，产生 $(-1)^{\deg\alpha}$，于是
\[
 d(\alpha\wedge\beta)=d\alpha\wedge\beta
 +(-1)^{\deg\alpha}\alpha\wedge d\beta.
\]
双线性再把结论推广到有限和。
\end{solution}

\begin{solution}{exterior-derivative-square-zero}
对 $\alpha=\sum_Ia_I dx^I$，
\[
 d^2\alpha=\sum_{I,j,k}\partial_k\partial_ja_I
 \,dx^k\wedge dx^j\wedge dx^I.
\]
$j=k$ 的项因楔积为零；$j\ne k$ 的两项系数由 Schwarz 定理相等，而
$dx^k\wedge dx^j=-dx^j\wedge dx^k$，故成对相消。前者使用交错性，后者
使用二阶混合偏导可交换。
\end{solution}

\begin{solution}{pullback-commutes-d}
对一次形式 $\omega=\sum_i a_i(y)dy^i$ 与 $F=(F^i)$，
$F^*\omega=\sum_i(a_i\circ F)dF^i$。于是
\[
 d(F^*\omega)=\sum_i d(a_i\circ F)\wedge dF^i
 =\sum_{i,j}(\partial_ja_i\circ F)dF^j\wedge dF^i=F^*(d\omega).
\]
函数情形是链式法则。一般形式写成函数与坐标一次形式的楔积，再由分次
Leibniz 公式归纳即可。
\end{solution}

\begin{solution}{stokes-rectangle-coordinate-proof}
设 $\omega=P\,dy\wedge dz+Q\,dz\wedge dx+R\,dx\wedge dy$。则
$d\omega=(P_x+Q_y+R_z)dx\wedge dy\wedge dz$。在长方体六个面上按外法向
诱导定向积分：$x$ 两面之和为 $\iiint P_x$，$y$ 两面之和为
$\iiint Q_y$，$z$ 两面之和为 $\iiint R_z$。三式来自一元基本定理与
Fubini，合并便是 $\int_{\partial Q}\omega=\int_Qd\omega$。
\end{solution}

\begin{solution}{boundary-orientation-square}
给平面区域取定向基 $(e_1,e_2)$。边界上“外法向优先”约定要求
$(n_{\rm out},\tau)$ 与 $(e_1,e_2)$ 同向。底边外法向 $-e_2$，故切向应为
$e_1$；右边依次向上，顶边向左，左边向下。所得方向正是逆时针方向。
\end{solution}

\begin{solution}{poincare-one-form-potential}
星形域星心取零，闭一次形式写成
$\omega=\sum_i a_i(x)dx^i$。同伦算子给出函数
\[
 \varphi(x)=\int_0^1\sum_i a_i(tx)x_i\,dt.
\]
求 $\partial_j$ 后利用闭性 $\partial_ja_i=\partial_i a_j$，被积式化为
$\frac d{dt}[t a_j(tx)]$，故 $\partial_j\varphi=a_j(x)$，即
$d\varphi=\omega$。
\end{solution}

\begin{solution}{poincare-homotopy-identity}
令 $F(t,x)=tx$，并以 $j_t(x)=(t,x)$ 表示切片嵌入。对 $k$-形式定义
\[
 H\omega=\int_0^1j_t^*\bigl(\iota_{\partial_t}F^*\omega\bigr)\,dt.
\]
先设 $k\ge1$。若
$\omega=\sum_Ia_I(x)\,dx^{i_1}\wedge\cdots\wedge dx^{i_k}$，则
$F^*dx^i=x_i\,dt+t\,dx^i$，从而
\[
 j_t^*\bigl(\iota_{\partial_t}F^*\omega\bigr)
 =\sum_Ia_I(tx)t^{k-1}
   \sum_{r=1}^k(-1)^{r-1}x_{i_r}\,
   dx^{i_1}\wedge\cdots\wedge\widehat{dx^{i_r}}
   \wedge\cdots\wedge dx^{i_k}.
\]
对这个有限和取 $d_x$。微分 $a_I(tx)$ 产生的各项，与
$j_t^*(\iota_{\partial_t}F^*d\omega)$ 中含偏导数的项逐项配对；微分
$x_{i_r}$ 产生的 $k$ 个项合为
$k t^{k-1}a_I(tx)\,dx^{i_1}\wedge\cdots\wedge dx^{i_k}$。连同
$t^k\frac d{dt}a_I(tx)$，正好得到
\[
 \frac d{dt}(F_t^*\omega)
 =d_xj_t^*\bigl(\iota_{\partial_t}F^*\omega\bigr)
  +j_t^*\bigl(\iota_{\partial_t}F^*d\omega\bigr).
\]
从 $0$ 到 $1$ 积分，左端为 $F_1^*\omega-F_0^*\omega$，右端为
$dH\omega+Hd\omega$。由于 $F_1=I$、$F_0=s_0$，得到
$dH+Hd=I-s_0^*$。当 $k=0$ 时，$H\omega=0$，而对函数使用一元微积分
基本定理可直接得到 $H(d\omega)=\omega-s_0^*\omega$。所有换序只涉及
有限和及连续函数的定积分。
\end{solution}

\begin{solution}{punctured-plane-closed-not-exact-form}
角形式 $\omega=(-y\,dx+x\,dy)/(x^2+y^2)$ 在
$\R^2\setminus\{0\}$ 上满足 $d\omega=0$。若它等于 $d\varphi$，则沿任何
闭曲线的积分应为零。但单位圆上 $\omega=dt$，故积分为 $2\pi$。矛盾说明
它闭而不恰当；障碍来自定义域的孔，不是局部微分计算。
\end{solution}

\begin{solution}{local-poincare-lemma}
给定开集 $U$ 中一点 $a$，可取足够小的球 $B(a,r)\subset U$。球关于 $a$
星形。若 $\omega$ 在 $U$ 上闭，则其限制在球上仍闭；星形域 Poincaré 引理
给出 $\omega|_{B(a,r)}=d\eta$。因此每个闭形式在每一点的某邻域上恰当。
\end{solution}

\begin{solution}{green-as-differential-form}
令 $\omega=Pdx+Qdy$，则
$d\omega=(Q_x-P_y)dx\wedge dy$。平面区域采用标准定向时，Stokes 公式
$\int_{\partial D}\omega=\int_Dd\omega$ 正好写成
\[
 \int_{\partial D}Pdx+Qdy=\iint_D(Q_x-P_y)dxdy,
\]
边界诱导定向即逆时针方向。
\end{solution}

\begin{solution}{gauss-as-differential-form}
对 $F=(P,Q,R)$ 取
$\omega=P\,dy\wedge dz+Q\,dz\wedge dx+R\,dx\wedge dy$。直接求外微分得
\[
 d\omega=(P_x+Q_y+R_z)dx\wedge dy\wedge dz
 =(\operatorname{div}F)dV.
\]
在定向曲面上，$\omega$ 的拉回等于 $F\cdot n\,dS$，所以 Stokes 公式恢复
散度定理。
\end{solution}

\begin{solution}{curl-as-exterior-derivative}
若 $\alpha=Pdx+Qdy+Rdz$，则
\[
 d\alpha=(R_y-Q_z)dy\wedge dz+(P_z-R_x)dz\wedge dx
 +(Q_x-P_y)dx\wedge dy.
\]
三个系数依次是 $\operatorname{curl}F$ 的三个分量。借助欧氏度量和定向把
二次形式识别为向量场后，$d\alpha$ 正对应旋度。
\end{solution}

\begin{solution}{forms-metric-independence-audit}
交错形式、楔积、外微分、拉回、定向以及 Stokes 公式只需光滑坐标结构和
定向，不需要长度或角度。把向量场变为一次形式、定义梯度、法向单位向量、
面积元以及把二次形式识别为旋度或通量形式，都使用欧氏内积；Hodge 星算子
更同时依赖度量与定向。因此微分形式版 Stokes 比向量分析的三个公式使用的
附加结构更少。
\end{solution}
